\section{System Design}


\subsection{Architectural Principles}

Crowdstart adheres to four design principles:

\textbf{Principle 1: Stateless Request Handling.} No request depends on state from a previous request to the same instance. All state resides in managed services.

\textbf{Principle 2: Idempotent Operations.} All mutating operations can be safely retried. This enables automatic retry on transient failures without application-level deduplication logic.

\textbf{Principle 3: Eventual Consistency by Default.} We accept eventual consistency for most operations, reserving strong consistency for the critical path (inventory decrements, payment capture).

\textbf{Principle 4: Event-Driven Processing.} Long-running operations execute asynchronously via task queues. Synchronous request handlers complete within milliseconds.

\subsection{Core Components}

The Crowdstart architecture comprises five primary components:

\subsubsection{API Gateway}

The gateway handles authentication, rate limiting, and request routing. Implemented as a stateless Go service, it validates JWT tokens against a cached public key and routes requests to appropriate backend services.

\begin{lstlisting}[language=Go,caption=Request routing]
func (g *Gateway) ServeHTTP(w http.ResponseWriter, r *http.Request) {
    ctx := appengine.NewContext(r)

    // Authenticate
    claims, err := g.auth.Validate(ctx, r)
    if err != nil {
        http.Error(w, "Unauthorized", 401)
        return
    }

    // Route to backend
    backend := g.router.Match(r.URL.Path)
    backend.ServeHTTP(w, r.WithContext(
        context.WithValue(ctx, "claims", claims),
    ))
}
\end{lstlisting}

\subsubsection{Product Catalog}

The catalog service manages product data: descriptions, pricing, images, and variants. Data resides in Google Cloud Datastore with Memcache providing read-through caching.

We model products as entities with multiple ancestor relationships:

\begin{lstlisting}[language=Go,caption=Product entity model]
type Product struct {
    ID          string    `datastore:"-"`
    StoreKey    *datastore.Key `datastore:"-"`
    Name        string
    Description string    `datastore:",noindex"`
    Price       int64     // cents
    Currency    string
    Variants    []Variant `datastore:"-"`
    CreatedAt   time.Time
    UpdatedAt   time.Time
}
\end{lstlisting}

Ancestor queries enable strongly consistent reads within a store's product hierarchy, while cross-store queries accept eventual consistency.

\subsubsection{Inventory Management}

Inventory presents the primary consistency challenge. We implement a reservation-based model:

\begin{enumerate}
\item Customer adds item to cart $\rightarrow$ soft reservation (TTL: 15 minutes).
\item Checkout initiated $\rightarrow$ hard reservation (transaction).
\item Payment confirmed $\rightarrow$ inventory decrement.
\item Payment failed $\rightarrow$ reservation released.
\end{enumerate}

Hard reservations use Datastore transactions with optimistic concurrency:

\begin{lstlisting}[language=Go,caption=Inventory reservation]
func (s *InventoryService) Reserve(ctx context.Context, sku string, qty int) error {
    key := datastore.NewKey(ctx, "Inventory", sku, 0, nil)

    _, err := datastore.RunInTransaction(ctx, func(tc context.Context) error {
        var inv Inventory
        if err := datastore.Get(tc, key, &inv); err != nil {
            return err
        }

        if inv.Available < qty {
            return ErrInsufficientInventory
        }

        inv.Available -= qty
        inv.Reserved += qty
        _, err := datastore.Put(tc, key, &inv)
        return err
    }, nil)

    return err
}
\end{lstlisting}

\subsubsection{Order Processing}

Orders transition through a state machine: \texttt{PENDING} $\rightarrow$ \texttt{PAID} $\rightarrow$ \texttt{FULFILLED} $\rightarrow$ \texttt{COMPLETED}. State transitions are driven by task queue handlers, ensuring durability and retry semantics.

\subsubsection{Payment Integration}

Payment processing integrates with Stripe via their Go library. We implement the payment flow as an idempotent operation using Stripe's idempotency keys:

\begin{lstlisting}[language=Go,caption=Idempotent payment capture]
func (p *PaymentService) Capture(ctx context.Context, order *Order) error {
    params := &stripe.ChargeParams{
        Amount:   stripe.Int64(order.Total),
        Currency: stripe.String(order.Currency),
        Source:   &stripe.SourceParams{Token: stripe.String(order.PaymentToken)},
    }
    params.SetIdempotencyKey(order.ID)

    _, err := charge.New(params)
    return err
}
\end{lstlisting}

\subsection{Data Model}

We employ a denormalized data model optimized for read-heavy workloads. Entity relationships are represented through key embedding rather than joins:

\begin{align}
\text{Store} &\rightarrow \text{Products} \rightarrow \text{Variants} \\
\text{Store} &\rightarrow \text{Orders} \rightarrow \text{LineItems} \\
\text{Customer} &\rightarrow \text{Addresses}, \text{PaymentMethods}
\end{align}

Denormalization introduces update anomalies. We mitigate this through:
\begin{itemize}
\item Immutable entities where possible (orders, transactions).
\item Background consistency jobs for derived data.
\item Version vectors for conflict resolution.
\end{itemize}

